\section{Hexad--octad incidence and the oriented functional}
\subsection{Dodecaphasic incidence and return functional}

The Barbero--Immirzi contribution uses angular transport and
hexad--octad incidence. Let $\mathsf H\subset\mathsf O$ be a hexad
and an octad with the hexad's marked pair preserved. The relevant flag
sets are
\begin{equation}
 \mathcal F_6=\mathsf H\times\binom{\mathsf H}{2}
 \lhook\joinrel\longrightarrow
 \mathcal F_8=\mathsf O\times\binom{\mathsf H}{2},
 \qquad |\mathcal F_6|=90,\quad |\mathcal F_8|=120.
 \label{iii:eq:barbero-flags}
\end{equation}
The cardinalities are $6\binom62$ and $8\binom62$; the same incidence
preserves the normalized defect
$(90-120)/(24\binom62)=-1/12$. This count identifies the two sectors
used by return, while the calendar determines their
relative multiplicities.

For $m\in\{90,120\}$, the tritic return on a contractive channel is
\begin{equation}
 S_m(q)=\sum_{j\geq0}q^{m+3mj}
       =\frac{q^m}{1-q^{3m}},\qquad 0<q<1.
 \label{iii:eq:barbero-return-series}
\end{equation}
Each term corresponds to the initial height $m$ and $j$ returns of
length $3m$. The geometric series thus preserves both the initial
height and return period. The dodecaphasic calendar
$\mathbb Z/12\mathbb Z$ supplies twelve sectors and a unique oriented
boundary incidence, denoted $\partial_{\rm ret}$. In the free
module of events, its fundamental chain is
\begin{equation}
 c_{12}^{+}
 =\sum_{j\in\mathbb Z/12\mathbb Z}[j,\mathcal F_6]
  -[\partial_{\rm ret},\mathcal F_8].
 \label{iii:eq:barbero-cycle}
\end{equation}
The linear readout $\mathscr L_q$ assigns $S_{90}(q)$ to each generator
$[j,\mathcal F_6]$ and $S_{120}(q)$ to the boundary generator. This yields
\begin{equation}
 \mathscr L_q(c_{12}^{+})=12S_{90}(q)-S_{120}(q).
 \label{iii:eq:barbero-weights}
\end{equation}
The weights $12$ and $-1$ are the chain's oriented multiplicities.
The conjugate orientation satisfies $c_{12}^{-}=-c_{12}^{+}$, with the
same incidence types.

\begin{proposition}[Singular coefficient of the fundamental return]
\label{iii:prop:barbero-pole}
The coefficient of $(1-q)^{-1}$ in
$\mathscr L_q(c_{12}^{+})$ is $1/24$.
\end{proposition}
\begin{proof}
For every integer $m>0$, the factorization
$1-q^{3m}=(1-q)\sum_{j=0}^{3m-1}q^j$ gives
\[
 \lim_{q\uparrow1}(1-q)S_m(q)=\frac1{3m}.
\]
Applying this identity to the functional already determined in
\eqref{iii:eq:barbero-weights} gives
\begin{equation}
 \lim_{q\uparrow1}(1-q)\mathscr L_q(c_{12}^{+})
 =\frac{12}{270}-\frac1{360}=\frac1{24}.
 \label{iii:eq:barbero-pole}
\end{equation}
In the complex coordinate $q-1$, the residue is $-1/24$.
\end{proof}

Singular evaluation verifies a consequence of incidence and tritic
return. The proof order is fixed by
\eqref{iii:eq:barbero-flags}--\eqref{iii:eq:barbero-weights}; the
singular coefficient is computed after the chain has been constructed.
